\documentclass[11pt]{article}

\usepackage[letterpaper,margin=1in]{geometry}
\usepackage{amsmath,amssymb}
\usepackage{booktabs}
\usepackage{longtable}
\usepackage{array}
\usepackage{xcolor}
\usepackage{listings}
\usepackage{url}
\usepackage{hyperref}
\usepackage{titlesec}
\usepackage{parskip}
\usepackage{fancyhdr}
\usepackage{enumitem}

\hypersetup{
  colorlinks=true,
  linkcolor=black,
  urlcolor=blue!50!black,
  citecolor=black,
  pdftitle={LonghornSilicon Cross-Block Architectural Conflict Register},
  pdfauthor={LonghornSilicon},
  pdfsubject={Architectural Conflict Register, Version conflicts-0.1}
}

\titleformat{\section}{\normalfont\Large\bfseries}{\thesection.}{0.6em}{}
\titleformat{\subsection}{\normalfont\large\bfseries}{\thesubsection}{0.6em}{}
\titleformat{\subsubsection}{\normalfont\normalsize\bfseries}{\thesubsubsection}{0.6em}{}

\pagestyle{fancy}
\fancyhf{}
\lhead{\small Cross-Block Conflict Register}
\rhead{\small \texttt{conflicts-0.1}}
\cfoot{\thepage}
\renewcommand{\headrulewidth}{0.4pt}

\setlist[itemize]{leftmargin=*,topsep=2pt,itemsep=1pt}

\title{\textbf{LonghornSilicon \\
Cross-Block Architectural Conflict Register}}
\author{LonghornSilicon}
\date{Version \texttt{conflicts-0.1} \quad 2026-05-22}

\begin{document}
\maketitle
\thispagestyle{fancy}

\begin{abstract}
\noindent
A standing register of design tensions between the four blocks of
the LonghornSilicon LLM inference accelerator: ACU (Attention
Compute Unit, block~1), KVCE (KV Cache Engine, block~2), TIU (Token
Importance Unit, block~3), and MHC (Memory Hierarchy Controller,
block~4). Each entry catalogues a specific incompatibility,
specification gap, or design choice that requires coordinated
resolution across more than one block. The register is structured
to be appended-to as integration work progresses, not rewritten
each iteration. The companion document
\textit{KVCE $\times$ ACU Integration Audit} (\texttt{kvce-acu-audit-0.2})
covers the empirical measurement of items that have been quantified;
this document is broader and catalogues items not yet measured.
\end{abstract}

\section{Scope and conventions}
\label{sec:scope}

The register covers conflicts and design tensions \emph{between}
LonghornSilicon blocks, not bugs within a single block (those are
filed as GitHub issues against the respective block's repository).
A conflict warrants an entry here when:

\begin{itemize}
\item It cannot be resolved in any single block's RTL or reference
model alone.
\item It requires a design decision that spans at least two of
\{ACU, KVCE, TIU, MHC, compiler\}.
\item It is not yet captured in a stable interface specification.
\end{itemize}

Each entry uses the following fields:

\begin{description}[leftmargin=2em,style=nextline]
\item[Severity] \textit{high} = blocks integrated correctness;
\textit{medium} = degrades quality or limits configuration but
chip runs; \textit{low} = inefficiency or design-overhead issue.
\item[Status] \textit{open} = needs decision; \textit{partially resolved} =
discussed and tracked but not concluded; \textit{resolved} = decision
made and implemented in some interface; \textit{deferred} = recognised
but parked beyond the current tape-out target.
\item[Affected blocks] which of \{ACU, KVCE, TIU, MHC, compiler\}
are touched.
\item[Resolution path] short description of how the conflict can be
closed and what the candidate options are.
\end{description}

\section{Summary table}

\begin{center}
\renewcommand{\arraystretch}{1.2}
\setlength{\tabcolsep}{4pt}
\begin{tabular}{@{}rp{0.32\textwidth}lll@{}}
\toprule
\textbf{\#} & \textbf{Conflict} & \textbf{Severity} & \textbf{Status} & \textbf{Affected} \\
\midrule
1 & Q4.12 input range vs.\ raw activation magnitudes
  & high   & open & ACU, KVCE \\
2 & FP16 routing depends on full-precision V available
  & high   & partially resolved & ACU, KVCE \\
3 & int16 Q4.12 $\leftrightarrow$ int8/fp16 bridge unspecified
  & high   & open & ACU, KVCE \\
4 & GQA convention not in KVCE ISA
  & medium & open & KVCE, compiler \\
5 & TIU cold-start eviction policy undecided
  & medium & open & TIU, ACU, MHC \\
6 & SRAM\_DEPTH default of 16 below realistic hot-tier need
  & medium & partially resolved & KVCE, MHC \\
7 & Eviction granularity vs.\ ACU tile size
  & medium & deferred & KVCE, MHC \\
8 & K (288b) vs.\ V (208b) word width waste
  & low    & deferred & KVCE \\
9 & eDRAM viability at 16FFC
  & medium & deferred & MHC, foundry \\
10 & Lloyd--Max centroids vs.\ early-layer activation distribution
  & medium & open & KVCE \\
11 & End-to-end accuracy on real LLM not measured
  & high   & open & ACU, KVCE, integration \\
12 & TurboQuant+ \texttt{turbo4} noise floor obviates PC FP16 path
  & informational & resolved & ACU, KVCE \\
\bottomrule
\end{tabular}
\end{center}

The remainder of this document expands each entry.

\section{Conflict 1 --- Q4.12 input range vs.\ raw activation magnitudes}
\label{sec:c1}

\textbf{Severity: high. Status: open. Affected: ACU, KVCE.}

KVCE accepts vectors as 64 lanes of int16 Q4.12 fixed-point, range
$\pm 8.0$ at resolution $2^{-12}$. Real Qwen2-0.5B post-\texttt{k\_proj}
key activations measured in this work have absolute maximum
approximately $\pm 152$ at the layers sampled. A naive
FP16~$\rightarrow$~Q4.12 conversion at the boundary clips
approximately $95\%$ of $K$'s magnitude information before any
compression takes place. The KVCE algorithm internally normalises
the vector by its L2 norm before quantization, so the clip is
unrecoverable: the post-clip vector has the wrong direction, and
no downstream step in KVCE can restore it.

\textbf{Resolution path.} Three candidates, none decided:

\begin{enumerate}
\item Widen the KVCE input format (e.g.\ Q8.8 with $\pm 128$ range,
or Q12.4 with $\pm 2048$). Breaks current RTL bit widths and all
existing reference-model tests.
\item Add a pre-normalisation stage between the ACU's
\texttt{k\_proj} output and the KVCE input port. The stage extracts
a per-vector scale factor (similar to KVCE's internal norm metadata
but at a wider range) and rescales the input into Q4.12-representable
range. Requires a new sub-block specification.
\item Document Q4.12 as the input contract and require upstream
software / firmware to normalise activations into this range before
streaming. Pushes the problem onto the compiler/runtime.
\end{enumerate}

\textbf{Status of tracking.} Filed as \texttt{kv-cache-engine\#2}
(see also \texttt{kvce\_acu\_integration\_audit.tex} \S5(a)).

\section{Conflict 2 --- FP16 routing depends on full-precision V}
\label{sec:c2}

\textbf{Severity: high. Status: partially resolved. Affected: ACU, KVCE.}

The precision controller's value proposition is per-tile FP16 routing
for ``peaked'' tiles where small attention-weight dimensions on
high-mass key positions are necessary to preserve output accuracy.
This argument was measured against $V$ supplied at FP16 (i.e.\ no
KV cache in the loop). When the chip integrates the KVCE, $V$ is
already lossy after $\sim$3-bit PolarQuant compression by the time it
reaches the ACU's FP16 SV matmul. The integration audit
(\texttt{kvce-acu-audit-0.2}) measures this empirically: the FP16
vs.\ INT8 distinction is $\sim 1{,}800\times$ smaller than the
KVCE-induced rMSE on the \texttt{turbo4} configuration.

The conflict is partially resolved in that the empirical answer is
in hand --- the PC's FP16 path provides no measurable benefit at
\texttt{turbo4}. It is not fully resolved because the design question
remains: do we (i)~retain the FP16 SV datapath knowing it is
under-utilised at \texttt{turbo4} but valuable at higher-precision
modes (\texttt{turbo8}, \texttt{turbo16}); (ii)~simplify the ACU to
INT8-only SV under the current configuration; (iii)~commit to a
higher-precision KVCE mode that restores the FP16 path's value?

\textbf{Resolution path.} Decision deferred until either the chip's
target KVCE compression mode is locked, or higher-precision KVCE
modes are implemented and re-measured. The ACU FP16 datapath
remains in the design pending that decision.

\section{Conflict 3 --- int16 Q4.12 $\leftrightarrow$ int8/fp16 bridge unspecified}
\label{sec:c3}

\textbf{Severity: high. Status: open. Affected: ACU, KVCE.}

KVCE decompresses each $K$ or $V$ vector to 64~lanes of int16 Q4.12.
The ACU's MAC array accepts either INT8 (for the INT8 path) or
fp32-at-fp16-precision (for the FP16 path). The conversion unit
that bridges these formats does not exist in any LonghornSilicon
specification, RTL, or reference model. It currently lives
implicitly inside ``kernel software''.

The bridge has at least two distinct sub-cases:

\begin{itemize}
\item Q4.12 $\rightarrow$ int8: per-tile symmetric quantisation
(\texttt{max\_abs} scan, divide by~127, round). Used on the INT8
SV path.
\item Q4.12 $\rightarrow$ float32-rounded-to-fp16: straightforward
multiply by $2^{-12}$ followed by IEEE-754 binary16 rounding. Used
on the FP16 path.
\end{itemize}

The bridge also has open questions about who picks the precision
(driven by the PC decision; routed by the bridge itself; or routed
upstream of the bridge); about where the per-tile scale lives (sent
alongside the decision; recomputed on every tile); and about
pipeline latency.

\textbf{Resolution path.} Specify the bridge as either (a)~an ACU
input-stage sub-block (preferred --- the precision-routing logic
naturally lives near the ACU), (b)~a KVCE output-stage sub-block
(less preferred --- KVCE shouldn't need to know about ACU's
INT8/FP16 routing), or (c)~its own block between ACU and KVCE
(over-architected for the role). Then write its RTL, reference
model, and ISA at the same level of formality as the existing block
specs.

\section{Conflict 4 --- GQA convention not in KVCE ISA}
\label{sec:c4}

\textbf{Severity: medium. Status: open. Affected: KVCE, compiler.}

Qwen2-class models use grouped-query attention (GQA): for example,
Qwen2-0.5B has 14~query heads and 2~key/value heads in a 7:1 ratio.
KVCE's ISA treats vectors as anonymous 64-D entities with no
\texttt{head} field on the AXI-Stream interface (only \texttt{tuser}
to select K vs.\ V path). The spec does not state whether:

\begin{itemize}
\item A separate KVCE instance is required per KV head, or
\item One KVCE instance multiplexes the KV heads through
address-space partitioning, or
\item The compiler is expected to flatten heads into the address
space, treating the block as head-agnostic.
\end{itemize}

\textbf{Resolution path.} Add a \texttt{KV\_HEAD} field to the ISA
register file, or document explicitly that the block is
head-agnostic and the compiler is responsible for head-indexed
addressing. The address-space approach is the cheapest in silicon
but ties the compiler tightly to the storage layout.

\section{Conflict 5 --- TIU cold-start eviction policy undecided}
\label{sec:c5}

\textbf{Severity: medium. Status: open. Affected: TIU, ACU, MHC.}

The Token Importance Unit (block~3) accumulates per-token attention
weight over decoding steps and uses the accumulator as the eviction
signal. In the first $\sim 50$--$100$ decode steps, accumulator
scores are roughly equal across tokens because the running sum has
not yet differentiated them. If the on-chip working set exceeds
SRAM capacity during this window, eviction by lowest-accumulator
is essentially random. This may cause the loss of tokens that
become important later (``heavy hitters'').

The conflict is not internal to the TIU: the eviction signal flows
to the KVCE (which advertises \texttt{evict\_needed}/\texttt{evict\_addr})
and the MHC (which spills/refills from off-chip DRAM). The choice
of cold-start policy affects all three.

\textbf{Resolution path.} Three published approaches available:
(i)~H2O's deferred-eviction-until-window-fills; (ii)~StreamingLLM's
pinned attention-sink tokens; (iii)~ForesightKV's learned-prior
pre-seeding of the accumulator. Option~(iii) is most principled and
slots into the TIU's planned interface without changing the
eviction comparator --- the accumulator is simply non-zero at the
start of decoding. ForesightKV adds a small offline-trained scorer
that runs at prefill; its hardware footprint and training pipeline
are open questions. See external inbound from A.~Singh (2026-05-20,
GitHub \texttt{abhi1432-boop/heavy-hitter-kv-cache}) for context.

\section{Conflict 6 --- \texttt{SRAM\_DEPTH=16} below hot-tier need}
\label{sec:c6}

\textbf{Severity: medium. Status: partially resolved. Affected: KVCE, MHC.}

The current KVCE \texttt{SRAM\_DEPTH} default is 16 entries, which
is the value that drives the OpenLane Sky130 signoff and area
numbers in the KVCE paper. For an LLM-inference workload at
seq\_len~$\geq 512$, the working set is much larger; even with
aggressive eviction the hot tier needs hundreds to thousands of
entries. The ISA range allows up to 4096, but no synthesis or
signoff has been run at deployment-realistic depths.

\textbf{Resolution path.} Re-run Yosys synth and LibreLane Sky130
signoff with \texttt{SRAM\_DEPTH=1024} (or whatever the target is)
to produce area and power numbers that reflect deployment.
This is recorded in the KVCE roadmap. The conflict is partially
resolved because the configurability exists; it remains open
because the signoff numbers we cite do not reflect a deployable
configuration.

\section{Conflict 7 --- Eviction granularity vs.\ ACU tile size}
\label{sec:c7}

\textbf{Severity: medium. Status: deferred. Affected: KVCE, MHC.}

KVCE evicts at vector granularity: one address, one
$K$~or~$V$~vector spilled. The ACU reads at $64 \times 64$ tile
granularity: 64 $K$~vectors and 64 $V$~vectors per attention tile.
If the MHC has to fetch evicted vectors from DRAM one at a time
on demand, the ACU's MAC pipeline stalls for many cycles per tile
on a cache miss. Burst fetch is required.

\textbf{Resolution path.} Constraint on the MHC: per-tile fetch
requests should be coalesced into single DRAM bursts (typically
$8$~or~$16$~vectors per burst at LPDDR4X). The TIU eviction policy
should evict at tile-aligned granularity when possible (a ``tile
eviction list'' rather than vector-by-vector). Detailed design
deferred to MHC specification.

\section{Conflict 8 --- K vs.\ V word-width waste}
\label{sec:c8}

\textbf{Severity: low. Status: deferred. Affected: KVCE.}

A compressed K entry is 288 bits; a compressed V entry is 208 bits.
\texttt{SRAM\_WIDTH} is set to $\max(288, 208) = 288$~bits, so each
V entry wastes 80~bits ($28\%$ of the word). The total SRAM area
overhead is meaningful at deployment-realistic depths but does not
affect functional correctness.

\textbf{Resolution path.} Either (i)~separate K and V SRAM macros
(doubles the macro count, more efficient bit utilisation), or
(ii)~packed allocation scheme (more complex address generation,
better bit utilisation), or (iii)~accept the overhead as the cost
of having a single uniform SRAM interface. Deferred to v0.2 of the
KVCE block.

\section{Conflict 9 --- eDRAM viability at 16FFC}
\label{sec:c9}

\textbf{Severity: medium. Status: deferred. Affected: MHC, foundry.}

The original chip block diagram describes the MHC as managing
``L1 SRAM (hottest KV) + L2 eDRAM 8--32~MB + off-chip LPDDR5''.
eDRAM at TSMC 16FFC is a specialty option, often requiring
additional mask layers and may not be included in standard university
program offerings (Muse Semi, TSMC University Program). If eDRAM is
unavailable, the L2 tier defaults to high-density SRAM with reduced
capacity per~mm$^2$.

\textbf{Resolution path.} Confirm with the foundry (TSMC University
Program directly, or via Muse Semi) which on-chip memory IPs are
available at 16FFC for student tape-outs. If eDRAM is unavailable,
revise the architecture to ``all-SRAM on-chip'' and accept the
capacity reduction; size LPDDR bandwidth accordingly. This decision
needs to land before MHC specification begins in earnest.

\section{Conflict 10 --- Lloyd-Max centroids vs.\ early-layer activations}
\label{sec:c10}

\textbf{Severity: medium. Status: open. Affected: KVCE.}

The KVCE compression's 8 Lloyd-Max centroids
($\pm 0.2451\sigma$, $\pm 0.7560\sigma$, $\pm 1.3439\sigma$,
$\pm 2.1520\sigma$ with $\sigma = 1/\sqrt{d}$) are derived under the
assumption that the Walsh--Hadamard-rotated, norm-normalised vector
has coordinates drawn from $\mathcal{N}(0, 1/\sqrt{d})$. This holds
for mid-network activations but breaks at layer~0 (and similar
near-input layers) because the model's residual stream has not yet
been smoothed by accumulated LayerNorms. The integration audit
measured a $16\times$ rMSE spread between layer~0 ($2.567$) and
layer~12 ($0.159$).

\textbf{Resolution path.} Candidate fixes are per-layer centroid
tables (small register file) or an adaptive
\texttt{turbo4-adaptive} mode in the compression algorithm. Filed
as \texttt{kv-cache-engine\#3}; see
\texttt{kvce\_acu\_integration\_audit.tex} \S5(b).

\section{Conflict 11 --- End-to-end accuracy not measured}
\label{sec:c11}

\textbf{Severity: high. Status: open. Affected: ACU, KVCE, integration.}

The integration audit measures per-tile relative MSE against a
dense FP16 attention baseline. The integrated rMSE of $0.36$ is
suggestive but does not directly bound model-output quality.
A serious chip-level accuracy claim requires measurement of
downstream metrics: perplexity on a held-out corpus,
HellaSwag accuracy, or similar, with the full attention pipeline
(ACU + KVCE round-trip + softmax + SV matmul) substituted for
dense FP16 attention. This has not been done.

\textbf{Resolution path.} Run a held-out evaluation
(perplexity on wikitext-2 or HellaSwag) against Qwen2-0.5B with
attention substituted by the integrated ACU $\times$ KVCE pipeline,
and compare to dense FP16. Filed as \texttt{kv-cache-engine\#4}.

\section{Conflict 12 --- \texttt{turbo4} noise floor obviates PC FP16 path}
\label{sec:c12}

\textbf{Severity: informational. Status: resolved (measured). Affected: ACU, KVCE.}

At the current KVCE configuration (\texttt{turbo4}: 3-bit PQ +
1-bit QJL for $K$, 3-bit PQ for $V$), the per-vector reconstruction
noise dominates the precision controller's INT8/FP16 routing
distinction. Quantified in
\texttt{kvce\_acu\_integration\_audit.tex} \S5(c): KV rMSE
$\sim 1{,}800\times$ the PC routing rMSE.

This is a design property of \texttt{turbo4}, not a defect. The
PC's FP16 path would re-acquire value at higher-precision modes
(\texttt{turbo8}, \texttt{turbo16}) where the per-vector noise
drops. Listed here for completeness.

\section{Closed items (for reference)}
\label{sec:closed}

\subsection*{KVCE reference-model reconstruction bug}

KVCE issue~\#1: \texttt{normalize()} used \texttt{coord\_frac=12}
instead of \texttt{norm\_frac=8} for left shift, causing $16\times$
overscale; C++ truncation division differed from Python
floor division. Fixed in \texttt{kv-cache-engine@9b1163a}; quality
gate added at \texttt{sw/reference\_model/test\_reconstruction\_quality.py}.
\textbf{Closed.}

\section*{Change log}

\begin{itemize}
\item \texttt{conflicts-0.1} (2026-05-22): initial register compiled
from the original cross-block conflict scan
(2026-05-15) plus residual items surfaced by the integration audit
(2026-05-19, 2026-05-22) plus TIU cold-start item from inbound
collaborator (2026-05-20).
\end{itemize}

\end{document}
